Zur \"Aquivalenz des Systems der Normalformen mit der Formenreihe $f$ bleibt jetzt nur noch zu zeigen, da\ss\ sich alle Formen der Formenreihe nach Polaren der Normalformen in dem unter 1. festgelegten Sinne entwickeln lassen. Es sei $\Pi$ ein Aggregat von Polaroperationen
\srcnumdisplay[4.2em]{(53.)}{%
 \left(\frac{\partial}{\partial\xi^{(i)}}\middle|\xi^{(k)}\right),
 \qquad (i\ne k;\ i=1,2,\ldots,\rho;\ k=1,2,\ldots,\rho)}
und $\Pi^\sigma$ ein Aggregat der speziellen Operationen (53.), f\"ur die $i=\sigma$ und $k<\sigma$. F\"ur eine Form $F$ der $\rho$ Reihen $\xi^{(1)},\ldots,\xi^{(\rho)}$, die in der Reihe $\xi^{(\rho)}$ vom Grade $k_\rho$ sei, gilt dann die Entwicklung\srcfn{*)}{F. Mertens, \emph{\"Uber eine Formel der Determinantentheorie}. Sitzungsber. d. Ak. d. Wiss., Wien, math.-naturw. Kl., Bd. 91. 2. Abt. (1885). Formel 20).}:
\srcnumdisplay[4.2em]{(54.)}{%
 F=\sum_{\alpha=0}^{k_\rho}\Pi F_\alpha,
 \qquad(\rho\leqq n)}
wo $F_\alpha$ in der letzten Reihe $\xi^{(\rho)}$ vom $\alpha$-ten Grade ist und aus $F$ durch die invarianten Differentialoperationen
\srcnumdisplay[4.2em]{(55.)}{%
 \left(\frac{\partial}{\partial\xi^{(1)}}\cdots\frac{\partial}{\partial\xi^{(\rho)}}\middle|\xi^{(1)}\cdots\xi^{(\rho)}\right)^\alpha}
und $\Pi^\rho$ hervorgeht. Ersetzt man bei der Erzeugung von $F_\alpha$ den Proze\ss\ (55.) durch
\srcnumdisplay[4.2em]{(56.)}{%
 \left(\frac{\partial}{\partial\xi^{(1)}}\cdots\frac{\partial}{\partial\xi^{(\rho)}}\middle|p_\rho\right)^\alpha,}
so entsteht eine Form $F_\alpha(p_\rho)$ mit nur $\rho-1$ Reihen $\xi^{(1)},\ldots,\xi^{(\rho-1)}$, auf die sich wieder (54.) anwenden l\"a\ss t. Die Fortsetzung des Verfahrens f\"uhrt zu Formen $G(p_\rho,p_{\rho-1},\ldots,p_1)$. Um aus diesen die Entwicklung von $F$ nach (54.) zu erhalten, hat man die einzelnen Reihen $p_\sigma$ durch $\xi^{(1)},\ldots,\xi^{(\sigma)}$ zu ersetzen und auf die entstehenden Formen den Polarproze\ss\ $\Pi$ (53.) anzuwenden; dies gibt (vgl. (48.)) eine Summe von transformierten Koeffizienten $(G)$. F\"ur $\rho=n$ bleibt beim ersten Schritt der Proze\ss\ (55.) erhalten. Bedeuten $c$ numerische Konstanten und wird noch abk\"urzend gesetzt
\srcnumdisplay[4.2em]{(57.)}{%
 (\xi^{(1)}\xi^{(2)}\cdots\xi^{(n)})=\Delta;
 \qquad
 \left(\frac{\partial}{\partial\xi^{(1)}}\frac{\partial}{\partial\xi^{(2)}}\cdots\frac{\partial}{\partial\xi^{(n)}}\right)=\nabla,}
so kommt im Falle $\rho=n$:
\srcnumdisplay[4.2em]{(58.)}{%
 F=\sum c\Delta^\alpha(G),}
w\"ahrend f\"ur $\rho<n$ nur $\alpha=0$, also $\sum c(G)$ rechts eintritt.
